
annulés par $\underline{I}$, donc de $\underline{O}_{X'}$-Modules, et on a :
\[
\cl(\underline{H})\in \operatorname{Im}K_\bullet(X'),\qquad
\cl(\underline{K})\in \operatorname{Im}K_\bullet(X'),
\]
d'où le résultat. On définit donc une fonction sur la catégorie des $\underline{O}_U$-Modules cohérents à valeurs dans $K_\bullet(X)/\operatorname{Im}K_\bullet(X')$ en associant à tout $\underline{O}_U$-Module cohérent la classe d'un de ses prolongements sur $X$.

Cette fonction est additive. En effet, soit :
\[
0\longrightarrow \underline{G}'\longrightarrow \underline{G}\longrightarrow \underline{G}''\longrightarrow 0
\]
une suite exacte de $\underline{O}_U$-Modules cohérents. Soit $\underline{F}$ un prolongement de $\underline{G}$ sur $X$; alors il existe un sous-Module $\underline{F}'$ de $\underline{F}$ prolongeant $\underline{G}'$, et $\underline{F}/\underline{F}'$ prolonge $\underline{G}''$. Elle se factorise donc par $K_\bullet(U)$, ce qui définit l'homomorphisme $K_\bullet(U)\to K_\bullet(X)/\operatorname{Im}K_\bullet(X')$ cherché. On vérifie immédiatement qu'il est inverse de $\rho$.

\textbf{Proposition 1.2.} \uline{Soient} $X,S$ \uline{deux schémas noethériens}, \uline{et} $f$ \uline{un morphisme} plat \uline{de} $X$ \uline{dans} $S$. \uline{On se donne une famille d'éléments} $(\xi_i)_{i\in I}$ \uline{de} $K^{\bullet}(X)$, \uline{et on considère} \uline{l'homomorphisme} :
\[
\varphi:(K_\bullet(S))^{(I)}\longrightarrow K_\bullet(X)
\]
\uline{défini par} :
\[
\varphi((x_i)_{i\in I})=\sum_{i\in I}\xi_i\cdot f^*(x_i),
\]
\uline{où} $f^*$ \uline{est l'homomorphisme} «~\uline{image inverse}~» \uline{de} $K_\bullet(S)$ \uline{dans} $K_\bullet(X)$, \uline{défini d'après} IV 2.12, \uline{et où} $K_\bullet(X)$ \uline{est considéré comme} $K^{\bullet}(X)$\uline{-module par} IV 2.10. \uline{On suppose que pour tout point} $s\in S$, \uline{les images des} $\xi_i$ \uline{dans} $K_\bullet(X_s)$ \uline{par} \uline{l'homomorphisme} :
\[
K^{\bullet}(X)\longrightarrow K^{\bullet}(X_s)\longrightarrow K_\bullet(X_s)
\]
\uline{engendrent} $K_\bullet(X_s)$ \uline{comme groupe abélien}. \uline{Alors} $\varphi$ \uline{est surjective}.

